\centerline{\Huge \bf  Олимпиада I}
\centerline{\Huge \bf 6 класс}

\begin{flushright}
    \scriptsize{Заключительный отбор в Сириус 2021г.\\ Кто загуглит - тому бан!}
\end{flushright}

\problem{1} 
Участок земли квадратной формы огорожен забором так, что столбы забора находятся в каждом углу и через равное расстояние по сторонам участка. Всего было использовано $24$ столба. Площадь земельного участка $144$ м$^{2}$. Найдите расстояние между двумя соседними столбами.

\problem{2}
В найденом археологами сундуке было $96,50$ рублей царскими монетами по 5 копеек, $10$ копеек, $25$ копеек, $50$ копеек и $1$ рублю. Если добавить в сундук одну монету $5$ копеек, две по 10 копеек, три по $25$ копеек, четыре по $50$ копеек и пять по $1$ рублю, то в сундуке будет одинаковое количество монет каждого вида. Подсчитайте, сколько монет каждого вида находится в сундуке.

\problem{3}
По кругу стоят $1000$ корзин. В каждой корзине, кроме одной, лежит по яблоку. За ход разрешается выбрать три соседних корзины $A, B, C$ ($B$ между $A$ и $C$) таких, что $A$ пуста, а в $B$ и $C$ - яблоки, и переложить яблоко из $C$ в $A$, съев яблоко из корзины $B$. Можно ли за несколько ходов съесть все яблоки кроме одного?

\problem{4}
Тест состоит из 6 вопросов, которые оцениваются в $1, 2, 3, 4, 5$ и $6$ баллов. Каждый вопрос оценивается в то количество баллов, каков его номер. Каждый ученик должен ответить на $6$ вопросов. При правильном ответе балл, присвоенный вопросу, прибавляется, а при неправильном ответе балл вычитается. Например, ученик, который правильно ответил только на вопросы $1, 3$ и $4,$ получает оценку, равную $1-2+3+4-5-6 = -5.$ Какое наибольшее количество учеников могло принимать участие в тесте, если оценки у всех различны?

\problem{5}
На красной тарелке лежат $20$ печений, а на синей $-$ $14$. Петя и Вася ходят по очереди, начинает Петя. Ход состоит в том, чтобы съесть два печенья, взятых с одной любой тарелки, или переместить одно печенье с красной тарелки на синюю (но не с синей тарелки на красную). Проигрывает тот, кто не может сделать ход. Кто из игроков может обеспечить себе победу независимо от игры противника?

\problem{6}
В таблице с $4$ строками и $k$ столбцами записаны натуральные числа так, что все числа в таблице различные и сумма чисел в каждом столбце равна $30$. Найдите максимально возможное количество столбцов в таблице.